% Fig. 2: where the collective units sit in the Vortex core.
\begin{tikzpicture}[x=1mm,y=1mm,font=\scriptsize,
  st/.style={draw,rounded corners=.6mm,minimum height=6mm,align=center,inner sep=1.5pt,fill=black!3},
  pe/.style={draw,minimum height=4.6mm,minimum width=17mm,align=center,inner sep=1.5pt},
  new/.style={pe,fill=corange!25,draw=corange!80!black,line width=.7pt},
  arr/.style={-{Latex[length=1.4mm]},line width=.6pt},
  lbl/.style={font=\tiny,inner sep=1pt}]
\node[st,minimum width=11mm] (fetch)  at (0,0)   {Fetch};
\node[st,minimum width=11mm] (decode) at (13,0)  {Decode};
\node[st,minimum width=14mm,minimum height=9mm] (issue) at (28,0) {Issue\\GPR 32 lanes};
\node[st,minimum width=11mm] (commit) at (71,0)  {Commit};
\draw[arr] (fetch) -- (decode);
\draw[arr] (decode) -- (issue);
% execute cluster
\node[pe] (alu) at (50,10.5) {ALU};
\node[pe] (mdv) at (50,5.2) {MulDiv};
\node[new] (ntt) at (50,-0.1) {NTT pair unit};
\node[new] (ksg) at (50,-5.4) {Keccak stage unit};
\node[new] (krd) at (50,-10.7) {Keccak round unit};
\node[draw=black!40,dashed,rounded corners=.8mm,fit=(alu)(krd),inner sep=1.6mm] (exe) {};
\node[lbl,anchor=south west] at (exe.north west) {execute};
\foreach \u in {alu,mdv,ntt,ksg,krd} {
  \draw[arr] (issue.east) -- ++(3,0) |- (\u.west);
  \draw[arr] (\u.east) -- ++(3,0) |- (commit.west);
}
\draw[arr] (commit.north) |- (60,17.5) -- node[lbl,above] {one destination per lane} (28,17.5) -- (issue.north);
% memory side
\node[st,minimum width=15mm] (lsu) at (28,-17) {LSU};
\node[st,minimum width=15mm] (mem) at (50,-17) {Cache / memory};
\node[st,minimum width=15mm,fill=black!10] (ptr) at (71,-17) {Pointer PE};
\draw[arr] (issue.south) -- (lsu.north);
\draw[arr,{Latex[length=1.4mm]}-{Latex[length=1.4mm]}] (lsu) -- (mem);
\draw[arr,{Latex[length=1.4mm]}-{Latex[length=1.4mm]}] (mem) -- node[lbl,below,yshift=-1pt] {200 B state} (ptr);
\node[lbl,anchor=north,text=black!60] at (50,-20.5) {memory-coupled baseline};
\end{tikzpicture}
